% !TeX program = xelatex
% MAU BAO CAO: bien dich bang XeLaTeX hoac LuaLaTeX, it nhat 2 lan.
% Quy tac chi tiet va prompt tai su dung: Huong_dan_va_prompt.txt.
% Day la mau thong nhat dua tren PDF da cung cap, khong phai quy dinh cua truong.
\documentclass[a4paper,12pt]{article}

% ===== 1. KHO GIAY, LE TRANG, FONT =====
\usepackage[left=3cm,right=2cm,top=2cm,bottom=2cm,footskip=1cm]{geometry}
\usepackage{fontspec}
% Uu tien font Times New Roman cai tren may hoac 4 tep TTF trong fonts/.
% Khong kem tep font thuong mai trong mau nay.
\IfFontExistsTF{Times New Roman}{
  \setmainfont{Times New Roman}
}{
  \IfFileExists{fonts/times.ttf}{
    \setmainfont{times.ttf}[
      Path=fonts/,
      BoldFont=timesbd.ttf,
      ItalicFont=timesi.ttf,
      BoldItalicFont=timesbi.ttf]
  }{
    \IfFontExistsTF{TeX Gyre Termes}{
      \setmainfont{TeX Gyre Termes}
      \PackageWarningNoLine{MauBaoCao}{Times New Roman unavailable; using TeX Gyre Termes}
    }{
      \setmainfont{DejaVu Serif}
      \PackageWarningNoLine{MauBaoCao}{Times New Roman unavailable; using DejaVu Serif}
    }
  }
}
\usepackage{indentfirst}
\usepackage{microtype}
\usepackage{graphicx}
\usepackage{tikz}
\usepackage{array,tabularx,booktabs,longtable}
\usepackage{enumitem}
\usepackage{titlesec}
\usepackage{tocloft}
\usepackage{caption}
\usepackage{xurl}
\usepackage[unicode,hidelinks]{hyperref}
\urlstyle{same}

% 13 pt la co chu; 19.5 pt la khoang cach giua hai duong chan chu.
% Dinh nghia truc tiep de khong nham voi he so cua goi setspace.
\renewcommand{\normalsize}{\fontsize{13}{19.5}\selectfont}
\AtBeginDocument{\normalsize}
\setlength{\parindent}{0.5cm}
\setlength{\parskip}{4pt}
\setlength{\emergencystretch}{2em}
\raggedbottom
\setlist{topsep=4pt,itemsep=2pt,parsep=0pt,leftmargin=*}

% ===== 2. BA CAP TIEU DE: 1 / 1.1 / 1.1.1 =====
\setcounter{secnumdepth}{3}
\setcounter{tocdepth}{3}
\renewcommand{\thesection}{\arabic{section}}
\renewcommand{\thesubsection}{\thesection.\arabic{subsection}}
\renewcommand{\thesubsubsection}{\thesubsection.\arabic{subsubsection}}
\titleformat{\section}[hang]
  {\normalfont\fontsize{16}{20}\selectfont\bfseries}
  {\thesection.}{0.6em}{}
\titleformat{\subsection}[hang]
  {\normalfont\fontsize{14}{18}\selectfont\bfseries}
  {\thesubsection}{0.6em}{}
\titleformat{\subsubsection}[hang]
  {\normalfont\fontsize{13}{17}\selectfont\bfseries}
  {\thesubsubsection}{0.6em}{}
\titlespacing*{\section}{0pt}{16pt}{8pt}
\titlespacing*{\subsection}{0pt}{12pt}{6pt}
\titlespacing*{\subsubsection}{0pt}{8pt}{4pt}

% ===== 3. MUC LUC VA SO TRANG =====
\renewcommand{\contentsname}{MỤC LỤC}
\renewcommand{\refname}{Tài liệu tham khảo}
\renewcommand{\cfttoctitlefont}{\fontsize{16}{20}\selectfont\bfseries}
\setlength{\cftbeforetoctitleskip}{0pt}
\setlength{\cftaftertoctitleskip}{18pt}
\renewcommand{\cftsecfont}{\normalfont\bfseries}
\renewcommand{\cftsecpagefont}{\normalfont\bfseries}
\renewcommand{\cftsubsecfont}{\normalfont}
\renewcommand{\cftsubsubsecfont}{\normalfont}
\renewcommand{\cftsecleader}{\cftdotfill{\cftdotsep}}
\cftsetindents{section}{0em}{2em}
\cftsetindents{subsection}{2em}{3em}
\cftsetindents{subsubsection}{5em}{4em}
\setlength{\cftbeforesecskip}{6pt}
\makeatletter
\def\ps@plain{%
  \let\@oddhead\@empty\let\@evenhead\@empty
  \def\@oddfoot{\hfil{\normalfont\fontsize{11}{13}\selectfont\thepage}\hfil}%
  \let\@evenfoot\@oddfoot}
\makeatother
\pagestyle{plain}

% ===== 4. BANG, HINH, CHU THICH =====
\renewcommand{\tablename}{Bảng}
\renewcommand{\figurename}{Hình}
\DeclareCaptionFont{reportcaption}{\fontsize{12}{15}\selectfont}
\captionsetup{font=reportcaption,labelfont=bf,labelsep=period,
  justification=centering,skip=6pt}
\newcolumntype{Y}{>{\raggedright\arraybackslash}X}
\newcommand{\bangfont}{\fontsize{12}{15}\selectfont}
% Bảng và hình đánh số toàn tài liệu: Bảng 1, Bảng 2; Hình 1, Hình 2.

% ===== 5. THONG TIN CAN THAY TREN BIA =====
\newcommand{\DonViChuQuan}{ĐẠI HỌC QUỐC GIA HÀ NỘI}
\newcommand{\TenTruong}{TRƯỜNG ĐẠI HỌC CÔNG NGHỆ}
\newcommand{\LoaiBaoCao}{REPORT PROGRESSION 1}
\newcommand{\TenMonHoc}{ĐIỆN TOÁN ĐÁM MÂY}
\newcommand{\TenDeTai}{SaTiBot - quản lý và tương tác với các tài nguyên AWS bằng ngôn ngữ tự nhiên}
\newcommand{\GiangVien}{Hoàng Văn Tùng}
\newcommand{\SoNhom}{10}
\newcommand{\LopMonHoc}{2627_INT3319_1}
\newcommand{\DiaDiemThoiGian}{Hà Nội, Tháng 10/2026}
\hypersetup{pdftitle={Mẫu báo cáo học phần},pdfauthor={Nhóm thực hiện}}

\begin{document}

% ===== 6. TRANG BIA =====
\begin{titlepage}
\thispagestyle{empty}
% Khung doi theo vung le: trai 3 cm, phai 2 cm, tren/duoi 2 cm.
\begin{tikzpicture}[remember picture,overlay]
  \draw[line width=2pt]
    ([xshift=3cm,yshift=-2cm]current page.north west)
    rectangle ([xshift=-2cm,yshift=2cm]current page.south east);
  \draw[line width=0.5pt]
    ([xshift=3.12cm,yshift=-2.12cm]current page.north west)
    rectangle ([xshift=-2.12cm,yshift=2.12cm]current page.south east);
\end{tikzpicture}\par
\centering
\setlength{\parindent}{0pt}
\setlength{\parskip}{0pt}
\vspace*{0.8cm}
{\fontsize{14}{18}\selectfont\bfseries\DonViChuQuan\par}
\vspace{0.15cm}
{\fontsize{15}{20}\selectfont\bfseries\TenTruong\par}
\vspace{0.25cm}
\rule{5cm}{0.4pt}\par
\vspace{0.5cm}
% Dat logo.png cung thu muc voi tep .tex; neu chua co se hien o giu cho.
\IfFileExists{logo.png}{
  \includegraphics[width=3.8cm,height=3.8cm,keepaspectratio]{logo.png}
}{
  \fbox{\parbox[c][3.4cm][c]{3.4cm}{\centering
    \fontsize{11}{14}\selectfont LOGO TRƯỜNG\\(logo.png)}}
}\par
\vspace{0.7cm}
{\fontsize{22}{27}\selectfont\bfseries\LoaiBaoCao\par}
\vspace{0.35cm}
\begin{minipage}{0.9\textwidth}\centering
  {\fontsize{15}{20}\selectfont\bfseries MÔN HỌC: \TenMonHoc\par}
\end{minipage}\par
\vspace{0.6cm}
{\fontsize{13}{17}\selectfont\bfseries ĐỀ TÀI\par}
\vspace{0.12cm}
\begin{minipage}{0.9\textwidth}\centering
  {\fontsize{16}{21}\selectfont\bfseries\TenDeTai\par}
\end{minipage}\par
\vspace{0.65cm}
{\fontsize{13}{19.5}\selectfont
\renewcommand{\arraystretch}{1.1}
\begin{tabular}{@{}>{\raggedright\arraybackslash\bfseries}p{4.8cm}>{\raggedright\arraybackslash}p{8cm}@{}}
  Giảng viên môn học: & \GiangVien\\
  Nhóm: & \SoNhom\\
  Sinh viên thực hiện: & Cao Anh Minh (Trưởng nhóm)\\
                      & Trịnh Thiên Quang\\
                      & Phạm Đức Hùng\\
                      & Đức Huy\\
  Lớp môn học: & \LopMonHoc
\end{tabular}\par}
\vfill
{\fontsize{13}{17}\selectfont\bfseries\DiaDiemThoiGian\par}
\vspace*{0.5cm}
\end{titlepage}

% ===== 7. MUC LUC TU DONG =====
\pagenumbering{roman}
\tableofcontents
\clearpage
\pagenumbering{arabic}

% ===== 8. NOI DUNG MAU: THAY BANG NOI DUNG BAO CAO =====
% Khong go so thu tu vao trong {Ten tieu de}.
\section{Đặt vấn đề}
Trình bày vấn đề nghiên cứu, lý do lựa chọn đề tài và ý nghĩa thực tiễn.
Đây là nội dung giữ chỗ để minh họa cách trình bày, cần thay bằng nội dung
thật của báo cáo trước khi nộp.

\subsection{Bối cảnh thực tế}
Giới thiệu bối cảnh, đối tượng sử dụng và nhu cầu cần giải quyết.

\subsubsection{Tình huống thứ nhất}
Mô tả tình huống, vấn đề phát sinh và nhu cầu hỗ trợ. Tiêu đề cấp 3
được đặt trên dòng riêng và xuất hiện trong mục lục.

\subsubsection{Tình huống thứ hai}
Mô tả một tình huống khác để làm rõ phạm vi của bài toán.

\subsection{Các khái niệm và phạm vi}
\begin{itemize}
  \item \textbf{Khái niệm chính:} Giải thích ngắn gọn thuật ngữ được sử dụng.
  \item \textbf{Phạm vi:} Nêu đối tượng, dữ liệu và giới hạn nghiên cứu.
\end{itemize}

\section{Khảo sát giải pháp và xác định khoảng trống}
\subsection{Tiêu chí đánh giá}
Nêu các tiêu chí dùng để đối chiếu những giải pháp liên quan.

\subsection{Các giải pháp liên quan}
Bảng~\ref{tab:khaosat} minh họa cách trình bày bảng. Trích dẫn tài liệu
theo số bằng lệnh trích dẫn tự động, ví dụ~\cite{nguonmau}.

\begin{table}[htbp]
\centering
\caption{Mẫu bảng khảo sát}
\label{tab:khaosat}
\bangfont
\begin{tabularx}{\linewidth}{@{}YYYY@{}}
\toprule
\textbf{Giải pháp} & \textbf{Tiêu chí 1} & \textbf{Tiêu chí 2} & \textbf{Tiêu chí 3}\\
\midrule
Giải pháp A & Nhận xét ngắn gọn. & Nhận xét ngắn gọn. & Nhận xét ngắn gọn.\\
Giải pháp B & Nhận xét ngắn gọn. & Nhận xét ngắn gọn. & Nhận xét ngắn gọn.\\
\bottomrule
\end{tabularx}
\end{table}

\subsection{Nhận xét và khoảng trống}
Tổng hợp kết quả khảo sát và xác định vấn đề còn cần nghiên cứu.

\section{Câu hỏi nghiên cứu}
\subsection{Câu hỏi thứ nhất}
Nêu câu hỏi nghiên cứu rõ ràng, gắn với vấn đề đã xác định.
\subsection{Câu hỏi thứ hai}
Nêu câu hỏi bổ sung nếu cần thiết.

\section{Mục tiêu dự kiến và giải pháp tổng quan}
\subsection{Mục tiêu dự kiến}
Nêu các mục tiêu cụ thể và kết quả đầu ra tương ứng.
\subsection{Giải pháp tổng quan}
Mô tả phương pháp hoặc hướng tiếp cận dự kiến.

% ===== 9. TAI LIEU THAM KHAO =====
\clearpage
\phantomsection
\addcontentsline{toc}{section}{Tài liệu tham khảo}
\begin{thebibliography}{99}
\bibitem{nguonmau}
[Tác giả]. [Năm]. \textit{[Tên tài liệu]}. [Nơi công bố].
[DOI hoặc URL thực tế, nếu có].
% Chi la cho giu mau, khong phai tai lieu tham khao thuc.
\end{thebibliography}

\end{document}
